\mychapter{ATZENI}

\section{Story of forensics}

The origins of forensics science can be interpreted in two ways:
\begin{itemize}
  \item as the moment when the term ``forensic science'' was formally defined
  \item as the earliest historical applications of forensic concepts aimed at identifying individuals and reconstructing events
\end{itemize}

\noindent
Some of the earliest forensic practices date back approximately 4,000 years.

\begin{itemize}
\item \textbf{Ancient Babylon}: fingerprints were used in commercial agreements to identify individuals involved in transactions.
\item \textbf{Ancient China}: the manuscript \emph{Xi Duan Hu} (``The Washing Away of Wrongs'') described early investigative techniques for determining causes of death, such as drowning or strangulation.
\end{itemize}


\subsubsection {Modern Forensic }

Modern forensic science began to emerge during the nineteenth century.\\
Key milestones include:

\begin{itemize}
\item \textbf{1835 Henry Goddard}: first documented use of ballistic comparison to identify the weapon used in a murder.
\item \textbf{Henry Faulds}: proposed systematic use of fingerprints for personal identification.
\item \textbf{Sir Francis Galton}: introduced structured fingerprint classification (loops, whorls, arches) and promoted large fingerprint collections.
\end{itemize}

\noindent
The scientific progress of forensic methods was largely driven by:
\begin{itemize}
\item methodological standardization,
\item structured data collection,
\item large datasets for comparison.
\end{itemize}


\subsubsection{Historical Evolution of Digital Forensics}

Digital forensics emerged much later than traditional forensic science, following the widespread adoption of computers and digital systems.
\bigskip

\noindent
A fundamental milestone in digital forensics is the \textbf{Morris Worm} of 1989, one of the first large-scale computer worms.
\medskip

\noindent
Its investigation required analysis of:
\begin{itemize}
\item network logs,
\item system activity,
\item email communications.
\end{itemize}

\noindent
Recent technological developments have significantly complicated digital investigations.

\begin{itemize}
  \item \textbf{Mobile devices} collect large amounts of personal information, including location data, communication history, and social interactions.
  \item \textbf{Cloud computing} introduces further complexity because data is frequently stored on remote servers rather than on the seized device. Investigators therefore often require cooperation from cloud providers or access
\end{itemize}


\subsection{The Morris Worm: Focus}

It is often considered both the first major cyber attack and one of the earliest cases that required a structured digital forensic investigation.\\
The worm was created by \textbf{Robert Morris}, a graduate student in computer science and the son of a cryptographer working at the U.S. National Security Agency. Morris did not intend to launch a destructive attack, but rather to experiment with the feasibility of a self-propagating worm.\\
In November 1988, the worm was released on the \textbf{arpanet}
\subsubsection*{Technical Characteristics}

The worm targeted \textbf{Unix systems} and was written in the C programming language.
\\
It propagated by exploiting several weaknesses common in systems of that period\\
However, Morris added a probability (approximately \textbf{14\%}) that the worm would still replicate even if the system reported that it was already infected. 

\subsubsection*{Impact}

The consequences were dramatic for the scale of the network at the time.

\begin{itemize}
\item Approximately \textbf{6,000 machines} were infected.
\item Around \textbf{2,000 infections occurred within 15 hours}.
\item Estimated economic damage reached \textbf{10 million USD}.
\end{itemize}

\subsubsection*{Legal Consequences}

Robert Tappan Morris was prosecuted under the \textbf{Computer Fraud and Abuse Act (CFAA)}, one of the first applications of this federal law.

The sentence included:
\begin{itemize}
\item three years of probation,
\item community service,
\item a \$10,000 fine.
\end{itemize}

\subsubsection*{Long-Term Effects}

The Morris Worm had several lasting impacts on cybersecurity:

\begin{itemize}
\item creation of the first \textbf{Computer Emergency Response Team (CERT)} at Carnegie Mellon University, supported by DARPA;
\item increased awareness of cybersecurity risks in academic networks;
\item recognition of the need for coordinated incident response.
\end{itemize}

\noindent
At the time, parts of the Internet backbone were disconnected for several days to prevent further propagation.
